\scanpage{151}
\mspage{V.23}
% archive sheet number: 150

\margin{(Cf. 5.0. sur les hypothèses préliminaires) \deleted{\unsure{Si le cas}}}
\uline{Corollaire} 5.3. \inserted{Supposons la condit. a) de 5.2. satisfaite.} La condition b) \inserted{de 5.2. (et donc la conclusion 5.2.1)} est satisfaite dans chacun des cas suivants (i) et (ii), où on a posé $d=\sup_{\deleted{\unsure{x'\in Z}}}\dim\overline{\{x'\}}$, où $x'$ parcourt les générisations de $x$ dans $X_0$ :

\deleted{\unsure{(i)}} \uline{Cas I} Pour toute générisation $x'$ de $x$ dans $X$, on a
\[
\left\{\begin{aligned}
\operatorname{prof\,ét}^{\unsure{*|\alpha}}_{x'}(X_0)&\geq m'+1&&\text{si }x'\in Z\\
\operatorname{prof\,ét}^{*}_{x'}(X_{\bar{\unsure{\eta|y}}})&\geq m'\deleted{\unsure{+1}}+\dim\overline{\{x\}}&&\text{si }x'\in X_\eta,
\end{aligned}\right.
\tag*{(5.3.1)}
\]
ce qui est le cas si on fait $m'=m-d-1$, avec $m$ un entier \deleted{$m$} tel que
\[
\left\{\begin{aligned}
\operatorname{prof\,ét}^{\unsure{\alpha|*}}_{x'}(X_0)&\geq m&&\text{pour toute \inserted{générisation} $x'\in Z$ de $x$}\\
\operatorname{prof\,ét}^{\alpha}_{x'}(X_{\bar{\eta}})&\geq m\deleted{\unsure{-d}}-1&&\text{pour tt gén. $x'$ de $x$ dans $X$, avec $x'\in X_\eta$.}
\end{aligned}\right.
\tag*{(5.3.2)}
\]

\uline{Cas II} \deleted{Pour toute générisation $x'$ de $x$ dans $Z$,} \inserted{\deleted{On a} Le morphisme $f$ est plat en $x$, et on a} \deleted{on a}
\[
m'=\operatorname{prof\,géom}^{*}_x(X_0)\deleted{\unsure{-d}},\quad\text{où}\quad
\operatorname{prof\,géom}^{*}_x(X_0)\overset{\mathrm{df}}{=}\operatorname{prof\,géom}_x(X_0)+\inserted{\dim\overline{\{x\}}}
\]
Cf. SGA 2 $\overline{\underline{\mathrm{XIV}}}$ 5.3 pour la notion de profondeur géométrique). \deleted{ce qui est le cas si on fait $m'=m-d$, avec $m$ un entier tel que}
\[
\deleted{\operatorname{prof\,géom}^{*}_x(X)}
\]

\uline{Démonstration}

\uline{Cas I} Les relations (5.3.1) étant stables par \inserted{générisation et} changement de $X$ en $X(\unsure{h|k})$, il suffit de vérifier qu'elles \deleted{\illegible{2}} entraînent
\[
\operatorname{prof\,ét}_x(X)\geq m'+1.
\tag*{(*)}
\]
Or la deuxième des relations (5.3.1) \deleted{\unsure{peut aussi}} \inserted{s'écrit} \deleted{s'écrire}
\[
\operatorname{prof\,ét}_{x'}(X')\deleted{\unsure{\geq m'}}+\dim\overline{\{x'\}}^{(\unsure{X_\eta|X})}\geq m'+\dim\overline{\{x\}},
\]
\margin{$=\dim\overline{\{x\}}+\dim\overline{\{\bar{x}'\}}^{(X')}$}
or
\[
\dim\overline{\{x'\}}^{(X_\eta)}=\dim\overline{\{x'\}}^{X}-1=\dim\overline{\{x\}}+\dim\overline{\{\bar{x}'\}}^{(X')}-1
\]
où $\bar{x}'$ est un pt de $X'$ au dessus de $x'$, donc on \unsure{peut|voit}
